% GENERATED FILE. Edit canonical lessons, then run npm run generate:books.
% canonical-source-hash: fnv1a64:842e68b70ff01931
% canonical-lessons: MR-C184-dhur, MR-C184-jyot, MR-C184-havaman, MR-C184-paryavaran, MR-C184-pradusan

\chapter{Smoke, Flame, Climate, Environment, Pollution}
\label{ch:mr-smoke-flame-climate-environment-pollution}
\hlchaptermodality{184}

\begin{chapteropening}
\textbf{By the end of this chapter:} \emph{I can say smoke, a flame, the climate, the environment and pollution.}

Complete the last lesson of chapter 184: I can say smoke, a flame, the climate, the environment and pollution.
\end{chapteropening}

\section[\texorpdfstring{\mr{धूर} (dhūr)}{धूर (dhūr)}]{\mr{धूर} (dhūr) — smoke}

\label{lesson:MR-C184-dhur}



\begin{quote}
Before the new one: say the Marathi for a monument, then the Marathi for a statue.
\end{quote}

\subsection*{You'll want to know: \mr{धूर}}
\textbf{\mr{धूर}} — \emph{dhūr} — \textquotedblleft{}smoke\textquotedblright{}.

\begin{grammarlens}[title={What you've built}]
One more word for this chapter.
\end{grammarlens}

\subsection*{Your turn}
\begin{itemize}
\raggedright
  \item \emph{Say it:} \emph{dhūr}
  \item \emph{Say it:} \emph{dhūr}, once more
  \item \emph{Say it:} say \emph{dhūr}
  \item \emph{Recall:} say the Marathi for a monument, then the Marathi for a statue, then say \emph{dhūr} again
\end{itemize}

\begin{tcolorbox}[breakable,colback=teal!4,colframe=teal!35!black,title={Before you move on}]
What does \mr{धूर} mean? (\textquotedblleft{}Smoke\textquotedblright{}.) Say it once more.
\end{tcolorbox}
\section[\texorpdfstring{\mr{ज्योत} (jyot)}{ज्योत (jyot)}]{\mr{ज्योत} (jyot) — a flame}

\label{lesson:MR-C184-jyot}



\begin{quote}
Before the new one: say the Marathi for a statue, then the Marathi for smoke.
\end{quote}

\subsection*{You'll want to know: \mr{ज्योत}}
\textbf{\mr{ज्योत}} — \emph{jyot} — \textquotedblleft{}a flame\textquotedblright{}.

\begin{grammarlens}[title={What you've built}]
One more word for this chapter.
\end{grammarlens}

\subsection*{Your turn}
\begin{itemize}
\raggedright
  \item \emph{Say it:} \emph{jyot}
  \item \emph{Say it:} \emph{jyot}, once more
  \item \emph{Say it:} say \emph{jyot}
  \item \emph{Recall:} say the Marathi for a statue, then the Marathi for smoke, then say \emph{jyot} again
\end{itemize}

\begin{tcolorbox}[breakable,colback=teal!4,colframe=teal!35!black,title={Before you move on}]
What does \mr{ज्योत} mean? (\textquotedblleft{}A flame\textquotedblright{}.) Say it once more.
\end{tcolorbox}
\section[\texorpdfstring{\mr{हवामान} (havāmān)}{हवामान (havāmān)}]{\mr{हवामान} (havāmān) — the climate, the weather}

\label{lesson:MR-C184-havaman}



\begin{quote}
Before the new one: say the Marathi for smoke, then the Marathi for a flame.
\end{quote}

\subsection*{You'll want to know: \mr{हवामान}}
\textbf{\mr{हवामान}} — \emph{havāmān} — \textquotedblleft{}the climate, the weather\textquotedblright{}.

\begin{grammarlens}[title={What you've built}]
One more word for this chapter.
\end{grammarlens}

\subsection*{Your turn}
\begin{itemize}
\raggedright
  \item \emph{Say it:} \emph{havāmān}
  \item \emph{Say it:} \emph{havāmān}, once more
  \item \emph{Say it:} say \emph{havāmān}
  \item \emph{Recall:} say the Marathi for smoke, then the Marathi for a flame, then say \emph{havāmān} again
\end{itemize}

\begin{tcolorbox}[breakable,colback=teal!4,colframe=teal!35!black,title={Before you move on}]
What does \mr{हवामान} mean? (\textquotedblleft{}The climate, the weather\textquotedblright{}.) Say it once more.
\end{tcolorbox}
\section[\texorpdfstring{\mr{पर्यावरण} (paryāvaraṇ)}{पर्यावरण (paryāvaraṇ)}]{\mr{पर्यावरण} (paryāvaraṇ) — the environment}

\label{lesson:MR-C184-paryavaran}



\begin{quote}
Before the new one: say the Marathi for a flame, then the Marathi for the climate.
\end{quote}

\subsection*{You'll want to know: \mr{पर्यावरण}}
\textbf{\mr{पर्यावरण}} — \emph{paryāvaraṇ} — \textquotedblleft{}the environment\textquotedblright{}.

\begin{grammarlens}[title={What you've built}]
One more word for this chapter.
\end{grammarlens}

\subsection*{Your turn}
\begin{itemize}
\raggedright
  \item \emph{Say it:} \emph{paryāvaraṇ}
  \item \emph{Say it:} \emph{paryāvaraṇ}, once more
  \item \emph{Say it:} say \emph{paryāvaraṇ}
  \item \emph{Recall:} say the Marathi for a flame, then the Marathi for the climate, then say \emph{paryāvaraṇ} again
\end{itemize}

\begin{tcolorbox}[breakable,colback=teal!4,colframe=teal!35!black,title={Before you move on}]
What does \mr{पर्यावरण} mean? (\textquotedblleft{}The environment\textquotedblright{}.) Say it once more.
\end{tcolorbox}
\section[\texorpdfstring{\mr{प्रदूषण} (pradūṣaṇ)}{प्रदूषण (pradūṣaṇ)}]{\mr{प्रदूषण} (pradūṣaṇ) — pollution}

\label{lesson:MR-C184-pradusan}



\begin{quote}
Before the new one: say the Marathi for the climate, then the Marathi for the environment.
\end{quote}

\subsection*{You'll want to know: \mr{प्रदूषण}}
\textbf{\mr{प्रदूषण}} — \emph{pradūṣaṇ} — \textquotedblleft{}pollution\textquotedblright{}.

\begin{grammarlens}[title={What you've built}]
One more word for this chapter.
\end{grammarlens}

\subsection*{Your turn}
\begin{itemize}
\raggedright
  \item \emph{Say it:} \emph{pradūṣaṇ}
  \item \emph{Say it:} \emph{pradūṣaṇ}, once more
  \item \emph{Say it:} say \emph{pradūṣaṇ}
  \item \emph{Recall:} say the Marathi for the climate, then the Marathi for the environment, then say \emph{pradūṣaṇ} again
\end{itemize}

\begin{tcolorbox}[breakable,colback=teal!4,colframe=teal!35!black,title={Before you move on}]
What does \mr{प्रदूषण} mean? (\textquotedblleft{}Pollution\textquotedblright{}.) Say it once more.
\end{tcolorbox}
